\section{System Architecture}
\label{sec:architecture}

This section describes the system end to end, from a user's Vietnamese message
to the grounded answer returned to the chat widget. The backend is a
Python/FastAPI service; the frontend is an existing React/Vite single-page
application with a floating chat widget. \Cref{fig:arch} sketches the data flow.

\begin{figure}[H]
\centering
\fbox{\begin{minipage}{0.92\textwidth}
\centering\small
\textbf{Frontend (React)} $\;\to\;$ \texttt{POST /api/chat} (text + optional image)\\[2pt]
$\big\downarrow$\\[2pt]
\textbf{RAG Pipeline:}\quad
Intent classify $\to$ [Pricing? $\Rightarrow$ Zalo redirect]\\
$\to$ Query enhancement (optional) $\to$ Retrieval (dense / BM25 / hybrid)\\
$\to$ Cross-encoder re-rank (optional) $\to$ LLM generation (grounded)\\[2pt]
$\big\downarrow$\\[2pt]
\textbf{Response:} answer + sources + redirect flag + metadata
\end{minipage}}
\caption{High-level request flow through the backend RAG pipeline.}
\label{fig:arch}
\end{figure}

\subsection{Offline stage: indexing}
The knowledge base lives in two JSON files (\texttt{products.json},
\texttt{business.json}) which are the single source of truth. At index time the
\emph{chunker} converts each product/document into one or more text chunks
(document-level chunking by default: one chunk per product, concatenating name,
ID, description, use cases, and technical specs). Each chunk carries metadata
(product/document \texttt{id}, category, name) --- the \texttt{id} is essential
for the retrieval metrics in \Cref{sec:methodology}. Chunks are embedded with a
Vietnamese sentence-transformer (\texttt{bkai-foundation-models/vietnamese-bi-encoder}, 768-dim)
and upserted into a persistent \textbf{ChromaDB} collection using cosine
distance. With 19 products + 5 company documents the corpus is 24 chunks.

\subsection{Online stage: the query pipeline}
For a text query, the orchestrator (\texttt{rag/pipeline.py}) executes the
following steps in order:

\begin{enumerate}[leftmargin=1.5em,itemsep=3pt]
  \item \textbf{Intent classification} (\texttt{rag/intent.py}). A cheap keyword
        heuristic catches unambiguous pricing and out-of-scope queries; anything
        else is sent to a single-label Gemini classifier that returns one of
        \{product, comparison, business, recommendation, pricing, image,
        out\_of\_scope\}.
  \item \textbf{Pricing short-circuit.} If the intent is \textit{pricing}, the
        pipeline returns a fixed Vietnamese message redirecting the user to Zalo
        \emph{without any retrieval or generation} --- this structurally
        guarantees the system never quotes a price.
  \item \textbf{Query enhancement} (optional, \texttt{rag/query\_enhancement.py}).
        For product/business/recommendation/comparison intents the query may be
        transformed by HyDE (generate a hypothetical product description and
        embed that \cite{gao2023hyde}), rewriting, or keyword expansion.
        \emph{This stage is disabled in the recommended configuration} for the
        reasons established in \Cref{sec:results}.
  \item \textbf{Retrieval} (\texttt{rag/retrieval.py}). One of three strategies:
        \emph{dense} (cosine search in ChromaDB), \emph{BM25}
        \cite{robertson2009bm25} over Vietnamese word-segmented tokens (via
        \texttt{pyvi}), or \emph{hybrid} --- the union of both candidate sets,
        with each retriever's scores min--max normalized to $[0,1]$ and combined
        as $\alpha \cdot \text{dense} + (1-\alpha)\cdot\text{BM25}$
        ($\alpha=0.5$).
  \item \textbf{Re-ranking} (optional, \texttt{rag/reranking.py}). A multilingual
        cross-encoder (\texttt{cross-encoder/mmarco-mMiniLMv2-L12-H384-v1})
        re-scores each (query, document) pair and keeps the top-$k$
        \cite{nogueira2019passage}. A cross-encoder jointly attends to the query
        and document, so it judges relevance far more precisely than the
        bi-encoder used for first-stage retrieval --- at the cost of one forward
        pass per candidate.
  \item \textbf{Generation} (\texttt{rag/generation.py}). The retrieved chunks
        are injected into a Vietnamese system prompt instructing the LLM to
        answer only from the provided context, never invent prices, cite the
        source product, and admit uncertainty. The default generator is Google
        \textit{Gemini 2.5 Flash}; an Ollama (OpenAI-compatible) backend is also
        supported for local open-source models.
\end{enumerate}

The pipeline returns a structured response: the answer text, the list of source
product/document names, the retrieved IDs, a \texttt{redirect\_to\_zalo} flag,
and a metadata block recording the configuration used (strategy, intent,
re-rank/enhancement flags) for reproducibility.

\subsection{Multimodal (image) path}
If the request includes an image, it is routed through a CLIP image matcher
(\texttt{open\_clip}, ViT-B/32, LAION-2B weights \cite{radford2021clip}). Product
reference images are embedded once and cached; the uploaded image is embedded
and matched by cosine similarity. If the best score is below a threshold the
system returns a polite Vietnamese apology; otherwise the top matches' documents
are retrieved and a description is generated. (This path is implemented but not
evaluated quantitatively in this report --- see \Cref{sec:discussion}.)

\subsection{Modularity for experiments}
Every component is selected through a single configuration object, and the
pipeline constructor accepts overrides for retrieval strategy, $\alpha$,
re-ranking, query-enhancement method, top-$k$, and LLM. This is what allows the
controlled, one-variable-at-a-time experiments in \Cref{sec:results}: an
experiment is just a YAML file specifying the configuration, and the runner
evaluates it against the fixed test set.
